\documentclass[12pt]{article}
\usepackage[a4paper,margin=1in]{geometry}
\usepackage[utf8]{inputenc}
\usepackage[T1]{fontenc}
\usepackage{amsmath,amssymb}
\usepackage{textcomp}
\usepackage{parskip}
\begin{document}
\begin{center}
{\Large \textbf{2013 STEP II}}
\end{center}

\subsection*{Page 1}
\begin{flushleft}
\bigskip\textbf{Section A: Pure Mathematics}\par
\bigskip\textbf{Question 1.} (i) Find the value of m for which the line y = mx touches the curve y = ln x .\par
\[If instead the line intersects the curve when x = a and x = b, where a < b, show that\]
\[ab = ba. Show by means of a sketch that a < e < b.\]
\textbf{(ii)} The line y = mx + c, where c > 0, intersects the curve y = ln x when x = p and x = q,\par
\[where p < q. Show by means of a sketch, or otherwise, that pq > qp .\]
\textbf{(iii)} Show by means of a sketch that the straight line through the points (p, ln p) and (q, ln q),\par
where e \leq p < q , intersects the y-axis at a positive value of y. Which is greater, \pie or\par
\[e\pi?\]
ln q ln p\par
\textbf{(iv)} Show, using a sketch or otherwise, that if 0 < p < q and - = e -1, then qp > pq.\par
q p\par
-\par
\bigskip\textbf{Question 2.} For n \geq 0, let\par
\[I = xn(1 x)ndx .\]
n\par
-\par
\[\leq\]
\textbf{(i)} For n \geq 1, show by means of a substitution that\par
1 1\par
\[xn-1 (1 x)ndx = xn(1 x)n-1 dx\]
- -\par
0 0\par
\[\leq \leq\]
and deduce that\par
\[2 xn-1 (1 x)ndx = I n-1 .\]
-\par
\[\leq\]
\[Show also, for n \geq 1, that\]
n\par
\[I =\]
xn-1\par
(1\par
x)n+1\par
dx\par
n\par
n + 1 -\par
\[\leq\]
n\par
\[and hence that I n = I n-1 .\]
2(2n + 1)\par
\textbf{(ii)} When n is a positive integer, show that\par
(n!)2\par
\[I = .\]
n\par
(2n + 1)!\par
\textbf{(iii)} Use the substitution x = sin 2 θ to show that I 1 = \pi 8 , and evaluate I 3 .\par
2 2\par
\end{flushleft}
\subsection*{Page 2}
\begin{flushleft}
\bigskip\textbf{Question 3.} (i) Given that the cubic equation x3 + 3ax2 + 3bx + c = 0 has three distinct real roots and\par
\[c < 0, show with the help of sketches that either exactly one of the roots is positive or\]
all three of the roots are positive.\par
\textbf{(ii)} Given that the equation x3 + 3ax2 + 3bx + c = 0 has three distinct real positive roots\par
show that\par
\[a > b > 0, a < 0, c < 0 . ( )\]
∗\par
[Hint: Consider the turning points.]\par
\textbf{(iii)} Given that the equation x3 + 3ax2 + 3bx + c = 0 has three distinct real roots and that\par
\[ab < 0, c > 0 ,\]
determine, with the help of sketches, the signs of the roots.\par
\textbf{(iv)} Show by means of an explicit example (giving values for a, b and c) that it is possible\par
for the conditions ( ) to be satisfied even though the corresponding cubic equation has\par
∗\par
only one real root.\par
\bigskip\textbf{Question 4.} The line passing through the point (a, 0) with gradient b intersects the circle of unit radius\par
centred at the origin at P and Q, and M is the midpoint of the chord P Q. Find the coordinates\par
of M in terms of a and b.\par
\textbf{(i)} Suppose b is fixed and positive. As a varies, M traces out a curve (the locus of M ).\par
\[Show that x = by on this curve. Given that a varies with 1 \leq a \leq 1, show that the\]
locus is a line se - gment of length 2b/(1 + b2) 2 1 . Give a sketch - showing the locus and the\par
unit circle.\par
\textbf{(ii)} Find the locus of M in the following cases, giving in each case its cartesian equation,\par
describing it geometrically and sketching it in relation to the unit circle:\par
\[(a) a is fixed with 0 < a < 1, and b varies with < b < ;\]
-\infty \infty\par
\[(b) ab = 1, and b varies with 0 < b \leq 1.\]
\end{flushleft}
\subsection*{Page 3}
\begin{flushleft}
\bigskip\textbf{Question 5.} (i) A function f(x) satisfies f(x) = f(1 x) for all x. Show, by differentiating with respect to\par
\[\leq 1 - 1 \leq\]
x, that f ( ) = 0 . If, in addition, f(x) = f( ) for all (non-zero) x, show that f ( 1) = 0\par
\[\leq 2 x -\]
\[and that f (2) = 0.\]
\textbf{(ii)} The function f is defined, for x = 0 and x = 1, by\par
( (\par
\[(x^2 x + 1)3\]
\[f(x) = - .\]
\[(x^2 x)2\]
-\par
\[Show that f(x) = f( ) and f(x) = f(1 x).\]
x -\par
\[Given that it has exactly three stationary points, sketch the curve y = f(x).\]
\textbf{(iii)} Hence, or otherwise, find all the roots of the equation f(x) = and state the ranges\par
\[of values of x for which f(x) > .\]
343\par
\[Find also all the roots of the equation f(x) = and state the ranges of values of x\]
343\par
\[for which f(x) > .\]
\bigskip\textbf{Question 6.} In this question, the following theorem may be used.\par
Let u1, u2, . . . be a sequence of (real) numbers. If the sequence is bounded above (that is,\par
u n \leq b for all n, where b is some fixed number) and increasing (that is, u n \geq u n-1 for all n),\par
then the sequence tends to a limit (that is, converges).\par
\[The sequence u^1, u^2, . . . is defined by u^1 = 1 and\]
\[u n+1 = 1 + (n \geq 1) . ( )\]
u ∗\par
n\par
\textbf{(i)} Show that, for n \geq 3,\par
u n u n-2\par
\[u n+2 u n = - .\]
- (1 + u\par
n\par
)(1 + u n-2)\par
\textbf{(ii)} Prove, by induction or otherwise, that 1 \leq u \leq 2 for all n.\par
n\par
\textbf{(iii)} Show that the sequence u1, u3, u5, . . . tends to a limit, and that the sequence u2, u4,\par
u6, . . . tends to a limit. Find these limits and deduce that the sequence u1, u2, u3, . . .\par
tends to a limit.\par
\[Would this conclusion change if the sequence were defined by ( ) and u^1 = 3?\]
∗\par
\end{flushleft}
\subsection*{Page 4}
\begin{flushleft}
7 (i) Write down a solution of the equation\par
2 2\par
\[x 2y = 1 , ( )\]
- ∗\par
for which x and y are non-negative integers.\par
\[Show that, if x = p, y = q is a solution of ( ), then so also is x = 3p + 4q, y = 2p + 3q.\]
∗\par
Hence find two solutions of ( ) for which x is a positive odd integer and y is a positive\par
∗\par
even integer.\par
\textbf{(ii)} Show that, if x is an odd integer and y is an even integer, ( ) can be written in the form\par
∗\par
2 1\par
\[n = m(m + 1) ,\]
where m and n are integers.\par
\textbf{(iii)} The positive integers a, b and c satisfy\par
3 4 2\par
\[b = c a ,\]
-\par
\[where b is a prime number. Express a and c^2 in terms of b in the two cases that arise.\]
Find a solution of a2 + b3 = c4, where a, b and c are positive integers but b is not prime.\par
\bigskip\textbf{Question 8.} The function f satisfies f(x) > 0 for x \geq 0 and is strictly decreasing (which means that\par
\[f(b) < f(a) for b > a).\]
\textbf{(i)} For t \geq 0, let A0(t) be the area of the largest rectangle with sides parallel to the\par
\[coordinate axes that can fit in the region bounded by the curve y = f(x), the y-axis and\]
\[the line y = f(t). Show that A^0(t) can be written in the form\]
\[A^0(t) = x^0 (f(x^0) f(t)) ,\]
-\par
\[\leq\]
\[where x^0 satisfies x0f (x^0) + f(x^0) = f(t) .\]
\textbf{(ii)} The function g is defined, for t > 0, by\par
1 t\par
\[g(t) = f(x)dx .\]
t\par
\[\leq\]
\[\leq\]
\[Show that tg (t) = f(t) g(t) .\]
-\par
\[Making use of a sketch show that, for t > 0,\]
t\par
\[(f(x) f(t)) dx > A^0(t)\]
-\par
\[\leq\]
\[and deduce that t2g \leq (t) > A^0(t).\]
-\par
\textbf{(iii)} In the case f(x) = , use the above to establish the inequality\par
1 + x\par
\[ln \sqrt\{1\} + t > 1 ,\]
\[- \sqrt\{1\} + t\]
\[for t > 0.\]
\end{flushleft}
\subsection*{Page 5}
\begin{flushleft}
\bigskip\textbf{Section B: Mechanics}\par
\bigskip\textbf{Question 9.} The diagram shows three identical discs in equilibrium in a vertical plane. Two discs rest, not\par
in contact with each other, on a horizontal surface and the third disc rests on the other two.\par
The angle at the upper vertex of the triangle joining the centres of the discs is 2θ.\par
θ\par
The weight of each disc is W . The coefficient of friction between a disc and the horizontal\par
surface is μ and the coefficient of friction between the discs is also μ.\par
\textbf{(i)} Show that the normal reaction between the horizontal surface and a disc in contact with\par
the surface is W .\par
\textbf{(ii)} Find the normal reaction between two discs in contact and show that the magnitude of\par
W sin θ\par
the frictional force between two discs in contact is .\par
2(1 + cos θ)\par
\textbf{(iii)} Show that if μ < 2 \sqrt3 there is no value of θ for which equilibrium is possible.\par
-\par
\bigskip\textbf{Question 10.} A particle is projected at an angle of elevation α (where α > 0) from a point A on horizontal\par
ground. At a general point in its trajectory the angle of elevation of the particle from A is θ\par
and its direction of motion is at an angle φ above the horizontal (with φ \geq 0 for the first half\par
\[of the trajectory and φ \leq 0 for the second half).\]
\[Let B denote the point on the trajectory at which θ = α and let C denote the point on the\]
\[trajectory at which φ = α.\]
- 2\par
\textbf{(i)} Show that, at a general point on the trajectory, 2 tan θ = tan α + tan φ .\par
\textasciicircum{}\circ\par
\textbf{(ii)} Show that, if B and C are the same point, then α = 60 .\par
\textasciicircum{}\circ\par
\textbf{(iii)} Given that α < 60 , determine whether the particle reaches the point B first or the\par
point C first.\par
\end{flushleft}
\subsection*{Page 6}
\begin{flushleft}
\bigskip\textbf{Question 11.} Three identical particles lie, not touching one another, in a straight line on a smooth horizontal\par
surface. One particle is projected with speed u directly towards the other two which are at\par
\[rest. The coefficient of restitution in all collisions is e, where 0 < e < 1 .\]
\textbf{(i)} Show that, after the second collision, the speeds of the particles are 1 u(1 e), 1 u(1 e2)\par
and 1 u(1 + e)2. Deduce that there will be a third collision whateve 2 r the - value 4 of e - .\par
\textbf{(ii)} Show that there will be a fourth collision if and only if e is less than a particular value\par
which you should determine.\par
\end{flushleft}
\subsection*{Page 7}
\begin{flushleft}
\bigskip\textbf{Section C: Probability and Statistics}\par
\bigskip\textbf{Question 12.} The random variable U has a Poisson distribution with parameter λ. The random variables\par
X and Y are defined as follows.\par
U if U is 1, 3, 5, 7, . . .\par
\[X =\]
\[\geq0 otherwise\]
U if U is 2, 4, 6, 8, . . .\par
\[Y =\]
\[\geq0 otherwise\]
\textbf{(i)} Find E(X) and E(Y ) in terms of λ, α and β, where\par
λ2 λ4 λ λ3 λ5\par
\[α = 1 + + + and β = + + + .\]
2! 4! \cdots 1! 3! 5! \cdots\par
\textbf{(ii)} Show that\par
λα + λ2β λ2α2\par
\[Var (X) =\]
α + β - (α + β)2\par
and obtain the corresponding expression for Var (Y ). Are there any non-zero values of λ\par
\[for which Var (X) + Var (Y ) = Var (X + Y ) ?\]
\bigskip\textbf{Question 13.} A biased coin has probability p of showing a head and probability q of showing a tail, where\par
\[p = 0, q = 0 and p = q. When the coin is tossed repeatedly, runs occur. A straight run of\]
( ( (\par
length n is a sequence of n consecutive heads or n consecutive tails. An alternating run of\par
length n is a sequence of length n alternating between heads and tails. An alternating run\par
can start with either a head or a tail.\par
Let S be the length of the longest straight run beginning with the first toss and let A be the\par
length of the longest alternating run beginning with the first toss.\par
\textbf{(i)} Explain why P(A = 1) = p2 + q2 and find P(S = 1). Show that P(S = 1) < P(A = 1).\par
\textbf{(ii)} Show that P(S = 2) = P(A = 2) and determine the relationship between P(S = 3) and\par
\[P(A = 3).\]
\textbf{(iii)} Show that, for n > 1, P(S = 2n) > P(A = 2n) and determine the corresponding\par
\[relationship between P(S = 2n + 1) and P(A = 2n + 1). [You are advised not to use\]
\[p + q = 1 in this part.]\]
\end{flushleft}
\end{document}
